\chapter{2004年全国I卷（文）}
\section{单选题}

\begin{problem}
设集合 \(U = \{1, 2, 3, 4, 5\}\)，\(A = \{1, 2, 3\}\)，\(B = \{2, 5\}\)，则 \(A \cap (\complement_U B) =\)（\quad）

\choices
    {\(\{2\}\)}
    {\(\{2, 3\}\)}
    {\(\{3\}\)}
    {\(\{1,3\}\)}
\end{problem}

\begin{answer}
D
\end{answer}

\begin{solution}
在全集 \(U\) 中，\(B\) 的补集为 \(\complement_U B=\{1,3,4\}\). 因此 \(A\cap(\complement_U B)=\{1,2,3\}\cap\{1,3,4\}=\{1,3\}\)，故选 D.
\end{solution}

\begin{problem}
已知函数 \(f(x)=\lg\dfrac{1-x}{1+x}\)，若 \(f(a)=\dfrac12\)，则 \(f(-a)=\)（\quad）

\choices
    {\(\dfrac12\)}
    {\(-\dfrac12\)}
    {\(2\)}
    {\(-2\)}
\end{problem}

\begin{answer}
B
\end{answer}

\begin{solution}
由对数真数为正，函数的定义域为 \(-1<x<1\). 对定义域内的任意 \(x\)，有 \(f(-x)=\lg\dfrac{1+x}{1-x}=\lg\left(\dfrac{1-x}{1+x}\right)^{-1}=-f(x)\). 因而 \(f(-a)=-f(a)=-\dfrac12\)，故选 B.
\end{solution}

\begin{problem}
已知 \(\bs{a}\)，\(\bs{b}\) 均为单位向量，它们的夹角为 \(60^{\circ}\)，那么 \(\left|\bs{a}+3\bs{b}\right|=\)（\quad）

\choices
    {\(\sqrt7\)}
    {\(\sqrt{10}\)}
    {\(\sqrt{13}\)}
    {\(4\)}
\end{problem}

\begin{answer}
C
\end{answer}

\begin{solution}
由数量积公式，\(\bs{a}\cdot\bs{b}=|\bs{a}|\,|\bs{b}|\cos60^{\circ}=\dfrac12\). 所以
\[
\left|\bs{a}+3\bs{b}\right|^2=|\bs{a}|^2+6\bs{a}\cdot\bs{b}+9|\bs{b}|^2=1+6\cdot\frac12+9=13
\]
模长非负，故 \(\left|\bs{a}+3\bs{b}\right|=\sqrt{13}\)，选 C.
\end{solution}

\begin{problem}
函数 \(y=\sqrt{x-1}+1\)（\(x\geqslant1\)）的反函数是（\quad）

\choices
    {\(y=x^{2}-2x+2\)（\(x<1\)）}
    {\(y=x^{2}-2x+2\)（\(x\geqslant1\)）}
    {\(y=x^{2}-2x\)（\(x<1\)）}
    {\(y=x^{2}-2x\)（\(x\geqslant1\)）}
\end{problem}

\begin{answer}
B
\end{answer}

\begin{solution}
原函数的定义域为 \([1,+\infty)\)，值域也是 \([1,+\infty)\). 交换自变量与因变量得 \(x=\sqrt{y-1}+1\)，由于 \(x\geqslant1\)，有 \(\sqrt{y-1}=x-1\). 两边平方可得 \(y-1=(x-1)^2\)，即 \(y=x^2-2x+2\)，且反函数的定义域为 \(x\geqslant1\)，故选 B.
\end{solution}

\begin{problem}
\(\left(2x^3-\dfrac1{\sqrt{x}}\right)^7\) 的展开式中常数项是（\quad）

\choices
    {\(14\)}
    {\(-14\)}
    {\(42\)}
    {\(-42\)}
\end{problem}

\begin{answer}
A
\end{answer}

\begin{solution}
展开式的第 \(k+1\) 项为 \(\mathrm C_7^k(2x^3)^{7-k}\left(-x^{-\frac12}\right)^k=\mathrm C_7^k2^{7-k}(-1)^kx^{21-\frac72k}\)，其中 \(k=0\)，\(1\)，\(\ldots\)，\(7\). 常数项要求 \(21-\dfrac72k=0\)，所以 \(k=6\). 此时常数项为 \(\mathrm C_7^6\cdot2\cdot(-1)^6=14\)，故选 A.
\end{solution}

\begin{problem}
设 \(\alpha\in\left(0,\dfrac\pi2\right)\)，若 \(\sin\alpha=\dfrac35\)，则 \(\sqrt2\cos\left(\alpha+\dfrac\pi4\right)=\)（\quad）

\choices
    {\(\dfrac75\)}
    {\(\dfrac15\)}
    {\(-\dfrac75\)}
    {\(-\dfrac15\)}
\end{problem}

\begin{answer}
B
\end{answer}

\begin{solution}
因为 \(\alpha\in\left(0,\dfrac\pi2\right)\)，所以 \(\cos\alpha=\sqrt{1-\sin^2\alpha}=\dfrac45\). 由两角和公式，\(\sqrt2\cos\left(\alpha+\dfrac\pi4\right)=\sqrt2\left(\cos\alpha\cos\dfrac\pi4-\sin\alpha\sin\dfrac\pi4\right)=\cos\alpha-\sin\alpha=\dfrac15\)，故选 B.
\end{solution}

\begin{problem}
椭圆 \(\dfrac{x^2}{4}+y^2=1\) 的两个焦点为 \(F_1\)，\(F_2\)，过 \(F_1\) 作垂直于 \(x\) 轴的直线与椭圆相交，一个交点为 \(P\)，则 \(\left|\overrightarrow{PF_2}\right|=\)（\quad）

\choices
    {\(\dfrac{\sqrt3}{2}\)}
    {\(\sqrt3\)}
    {\(\dfrac72\)}
    {\(4\)}
\end{problem}

\begin{answer}
C
\end{answer}

\begin{solution}
椭圆的长半轴为 \(a=2\)，短半轴为 \(b=1\)，所以焦半距 \(c=\sqrt{a^2-b^2}=\sqrt3\). 过焦点作的直线垂直于 \(x\) 轴，故交点 \(P\) 与 \(F_1\) 的横坐标相同，从而 \(x_P^2=3\). 代入椭圆方程得 \(y_P^2=1-\dfrac34=\dfrac14\)，从而 \(|PF_1|=\dfrac12\). 椭圆上任一点满足 \(|PF_1|+|PF_2|=2a=4\)，所以 \(|PF_2|=4-\dfrac12=\dfrac72\)，故选 C.
\end{solution}

\begin{problem}
设抛物线 \(y^2=8x\) 的准线与 \(x\) 轴交于点 \(Q\)，若过点 \(Q\) 的直线 \(l\) 与抛物线有公共点，则直线 \(l\) 的斜率的取值范围是（\quad）

\choices
    {\(\left[-\dfrac12,\dfrac12\right]\)}
    {\([-2,2]\)}
    {\([-1,1]\)}
    {\([-4,4]\)}
\end{problem}

\begin{answer}
C
\end{answer}

\begin{solution}
由 \(y^2=8x=4\cdot2x\)，抛物线的准线为 \(x=-2\)，故 \(Q=(-2,0)\). 先取斜率 \(k\ne0\)，则直线 \(l\) 为 \(y=k(x+2)\). 与抛物线联立得 \(k^2(x+2)^2=8x\)，即 \(k^2x^2+(4k^2-8)x+4k^2=0\). 该方程有实根的条件为
\[
\Delta=(4k^2-8)^2-16k^4=64(1-k^2)\geqslant0
\]
所以 \(|k|\leqslant1\). 当 \(0<k^2\leqslant1\) 时，二次方程两根之和为 \(\dfrac{8-4k^2}{k^2}\geqslant0\)，两根之积为 \(4>0\)，故两根均为正数，确实对应抛物线上的交点；当 \(k=0\) 时，直线 \(y=0\) 经过抛物线顶点，也满足条件. 因此斜率范围为 \([-1,1]\)，故选 C.
\end{solution}

\begin{problem}
为了得到函数 \(y=\sin\left(2x-\dfrac\pi6\right)\) 的图象，可以将函数 \(y=\cos2x\) 的图象（\quad）

\choices
    {向右平移 \(\dfrac\pi6\) 个单位长度}
    {向右平移 \(\dfrac\pi3\) 个单位长度}
    {向左平移 \(\dfrac\pi6\) 个单位长度}
    {向左平移 \(\dfrac\pi3\) 个单位长度}
\end{problem}

\begin{answer}
B
\end{answer}

\begin{solution}
将 \(y=\cos2x\) 向右平移 \(h\) 个单位后，函数变为 \(y=\cos2(x-h)=\sin\left(2x-2h+\dfrac\pi2\right)\). 要与 \(\sin\left(2x-\dfrac\pi6\right)\) 相同，应满足 \(\dfrac\pi2-2h=-\dfrac\pi6\)，解得 \(h=\dfrac\pi3\)，故选 B.
\end{solution}

\begin{problem}
已知正四面体 \(ABCD\) 的表面积为 \(S\)，其四个面的中心分别为 \(E\)，\(F\)，\(G\)，\(H\). 设四面体 \(EFGH\) 的表面积为 \(T\)，则 \(\dfrac TS\) 等于（\quad）

\choices
    {\(\dfrac19\)}
    {\(\dfrac49\)}
    {\(\dfrac14\)}
    {\(\dfrac13\)}
\end{problem}

\begin{answer}
A
\end{answer}

\begin{solution}
设正四面体的棱长为 \(s\)，并以其重心为原点，记四个顶点的位置向量为 \(\bs{a}\)，\(\bs{b}\)，\(\bs{c}\)，\(\bs{d}\). 由于重心在原点，\(\bs{a}+\bs{b}+\bs{c}+\bs{d}=\bs0\). 设面 \(BCD\)，\(ACD\)，\(ABD\)，\(ABC\) 的中心依次为 \(E\)，\(F\)，\(G\)，\(H\)，则 \(\overrightarrow{OE}=-\dfrac13\bs{a}\)，\(\overrightarrow{OF}=-\dfrac13\bs{b}\)，\(\overrightarrow{OG}=-\dfrac13\bs{c}\)，\(\overrightarrow{OH}=-\dfrac13\bs{d}\). 因而 \(EF=\dfrac13|\bs{a}-\bs{b}|=\dfrac s3\)，\(EG=\dfrac13|\bs{a}-\bs{c}|=\dfrac s3\)，\(EH=\dfrac13|\bs{a}-\bs{d}|=\dfrac s3\)，\(FG=\dfrac13|\bs{b}-\bs{c}|=\dfrac s3\)，\(FH=\dfrac13|\bs{b}-\bs{d}|=\dfrac s3\)，\(GH=\dfrac13|\bs{c}-\bs{d}|=\dfrac s3\). 所以 \(EFGH\) 是棱长为 \(s/3\) 的正四面体. 面积按长度平方缩放，故 \(T/S=(1/3)^2=1/9\)，选 A.
\end{solution}

\begin{problem}
从数字 \(1\)，\(2\)，\(\cdots\)，\(9\) 中，随机抽取 \(3\) 个不同的数，则这 \(3\) 个数的和为偶数的概率为（\quad）

\choices
    {\(\dfrac59\)}
    {\(\dfrac49\)}
    {\(\dfrac{11}{21}\)}
    {\(\dfrac{10}{21}\)}
\end{problem}

\begin{answer}
C
\end{answer}

\begin{solution}
从 \(1\) 到 \(9\) 有 \(5\) 个奇数和 \(4\) 个偶数，等可能的取法共有 \(\mathrm C_9^3\) 种. 三个数的和为偶数，当且仅当取到三个偶数，或取到两个奇数和一个偶数. 有利取法数为 \(\mathrm C_4^3+\mathrm C_5^2\mathrm C_4^1=4+40=44\)，所以所求概率为 \(\dfrac{44}{\mathrm C_9^3}=\dfrac{44}{84}=\dfrac{11}{21}\)，故选 C.
\end{solution}

\begin{problem}
\(a^2+b^2=1\)，\(b^2+c^2=2\)，\(c^2+a^2=2\)，\(ab+bc+ca\) 的最小值为（\quad）

\choices
    {\(\sqrt3-\dfrac12\)}
    {\(\dfrac12-\sqrt3\)}
    {\(-\dfrac12-\sqrt3\)}
    {\(\dfrac12+\sqrt3\)}
\end{problem}

\begin{answer}
B
\end{answer}

\begin{solution}
第二个等式减去第三个等式，得 \(b^2-a^2=0\)，所以 \(a^2=b^2\). 结合 \(a^2+b^2=1\)，得 \(a^2=b^2=\dfrac12\)，再由任一其余等式得 \(c^2=\dfrac32\). 因此 \(a\)，\(b=\pm\dfrac1{\sqrt2}\)，\(c=\pm\sqrt{\dfrac32}\).

当 \(a\)，\(b\) 异号时，\(ab=-\dfrac12\)，且 \(a+b=0\)，故 \(ab+bc+ca=-\dfrac12\). 当 \(a\)，\(b\) 同号时，\(ab=\dfrac12\)，且 \(|a+b|=\sqrt2\). 为使 \(c(a+b)\) 最小，应取 \(c\) 与 \(a+b\) 异号，此时 \(c(a+b)=-\sqrt{\dfrac32}\sqrt2=-\sqrt3\)，从而 \(ab+bc+ca=\dfrac12-\sqrt3\). 因为 \(\dfrac12-\sqrt3< -\dfrac12\)，最小值为 \(\dfrac12-\sqrt3\)，故选 B.
\end{solution}

\section{填空题}

\begin{problem}
不等式 \(x+x^3\geqslant0\) 的解集是\fillinblank{}.
\end{problem}

\begin{answer}
\(\{x\mid x\geqslant0\}\)
\end{answer}

\begin{solution}
因式分解得 \(x+x^3=x(1+x^2)\). 对任意实数 \(x\)，都有 \(1+x^2>0\)，所以原不等式等价于 \(x\geqslant0\). 解集为 \(\{x\mid x\geqslant0\}\).
\end{solution}

\begin{problem}
已知等比数列 \(\{a_n\}\) 中，\(a_3=3\)，\(a_{10}=384\)，则该数列的通项 \(a_n=\)\fillinblank{}.
\end{problem}

\begin{answer}
\(3\cdot2^{n-3}\)
\end{answer}

\begin{solution}
设等比数列的公比为 \(q\). 由等比数列通项关系，\(a_{10}=a_3q^7\)，所以 \(3q^7=384=3\cdot2^7\)，即 \(q^7=2^7\). 由于七次方函数在实数集上严格递增，得 \(q=2\). 因而 \(a_n=a_3q^{n-3}=3\cdot2^{n-3}\).
\end{solution}

\begin{problem}
由动点 \(P\) 向圆 \(x^2+y^2=1\) 引两条切线 \(PA\)，\(PB\)，切点分别为 \(A\)，\(B\)，\(\angle APB=60^{\circ}\)，则动点 \(P\) 的轨迹方程为\fillinblank{}.
\end{problem}

\begin{answer}
\(x^2+y^2=4\)
\end{answer}

\begin{solution}
设圆心为 \(O\). 因为半径垂直切线，所以 \(OA\perp PA\)，\(OB\perp PB\)，四边形 \(OAPB\) 在 \(A\)，\(B\) 处的两个角都是直角，从而 \(\angle AOB+\angle APB=180^{\circ}\). 由 \(\angle APB=60^{\circ}\)，得 \(\angle AOB=120^{\circ}\).

又因为从同一点引圆的两条切线长度相等，\(PA=PB\)，同时 \(OA=OB\)，所以 \(OP\) 是 \(\angle AOB\) 的角平分线. 在直角三角形 \(OAP\) 中，\(\angle AOP=60^{\circ}\)，且 \(OA=1\)，故 \(OP=\dfrac{OA}{\cos60^{\circ}}=2\). 因此点 \(P(x,y)\) 满足 \(x^2+y^2=OP^2=4\). 反过来，圆 \(x^2+y^2=4\) 上任一点到 \(O\) 的距离为 \(2>1\)，均可引出两条切线，且由上述关系切线夹角为 \(60^{\circ}\)，所以轨迹方程为 \(x^2+y^2=4\).
\end{solution}

\begin{problem}
已知 \(a\)，\(b\) 为不垂直的异面直线，\(\alpha\) 是一个平面，则 \(a\)，\(b\) 在 \(\alpha\) 上的射影有可能是：

\begin{circlelist}
\item 两条平行直线；
\item 两条互相垂直的直线；
\item 同一条直线；
\item 一条直线及其外一点.
\end{circlelist}

在上面结论中，正确结论的编号是\fillinblank{}.
\end{problem}

\begin{answer}
\(\circled{1}\)、\(\circled{2}\)、\(\circled{4}\)
\end{answer}

\begin{solution}
取投影平面为 \(z=0\)，并把正投影方向取为 \(z\) 轴方向. 例如取 \(a:(x,y,z)=(t,0,t)\)，\(b:(x,y,z)=(s,1,2s)\). 两直线方向向量分别为 \((1,0,1)\)，\((1,0,2)\)，不平行且内积为 \(3\ne0\)，两直线又因第二个坐标不同而不相交，所以它们是异面且不垂直的；它们的投影分别为 \(y=0\) 与 \(y=1\) 两条平行直线，故结论 \(\circled{1}\) 可能成立.

取 \(a:(x,y,z)=(t,0,t)\)，\(b:(x,y,z)=(0,s,s+1)\). 两直线方向向量分别为 \((1,0,1)\)，\((0,1,1)\)，内积为 \(1\ne0\)，且联立三坐标无解，故两直线异面且不垂直；它们的投影为 \(x\) 轴和 \(y\) 轴，故结论 \(\circled{2}\) 可能成立.

若两条原直线的投影是同一条直线 \(m\)，则每条原直线都包含在由 \(m\) 和投影方向所确定的同一个平面内. 两条共面直线不可能是异面直线，所以结论 \(\circled{3}\) 不成立.

取 \(a:(x,y,z)=(t,0,t)\)，\(b:(x,y,z)=(1,1,s)\). 它们的方向向量内积为 \(1\ne0\)，且两直线不相交、不平行，所以是异面且不垂直的；投影分别为直线 \(y=0\) 和点 \((1,1)\)，该点在直线 \(y=0\) 外，故结论 \(\circled{4}\) 可能成立. 因此正确结论为 \(\circled{1}\)、\(\circled{2}\)、\(\circled{4}\).
\end{solution}

\section{解答题}

\begin{problem}
等差数列 \(\{a_n\}\) 的前 \(n\) 项和记为 \(S_n\)，已知 \(a_{10}=30\)，\(a_{20}=50\).

\begin{enumerate}
\item 求通项 \(a_n\)；
\item 若 \(S_n=242\)，求 \(n\).
\end{enumerate}
\end{problem}

\begin{solution}
\begin{enumerate}
\item 设首项为 \(a_1\)，公差为 \(d\). 由等差数列通项公式，\(a_{10}=a_1+9d=30\)，\(a_{20}=a_1+19d=50\). 两式相减得 \(10d=20\)，所以 \(d=2\)，代回得 \(a_1=12\). 因而 \(a_n=a_1+(n-1)d=12+2(n-1)=2n+10\).
\item 由等差数列前 \(n\) 项和公式，\(S_n=\dfrac{n(a_1+a_n)}2=\dfrac{n(12+2n+10)}2=n(n+11)\). 由 \(S_n=242\)，得 \(n^2+11n-242=0\)，即 \((n-11)(n+22)=0\). 因为项数 \(n\) 是正整数，所以 \(n=11\).
\end{enumerate}
\end{solution}

\begin{problem}
求函数 \(f(x)=\dfrac{\sin^4x+\cos^4x+\sin^2x\cos^2x}{2-\sin2x}\) 的最小正周期、最大值和最小值.
\end{problem}

\begin{solution}
令 \(s=\sin2x\). 由 \(\sin^2x\cos^2x=\dfrac14\sin^22x=\dfrac14s^2\)，以及 \(\sin^4x+\cos^4x=(\sin^2x+\cos^2x)^2-2\sin^2x\cos^2x=1-\dfrac12s^2\)，得分子为 \(1-\dfrac14s^2=\dfrac{(2-s)(2+s)}4\). 因为 \(-1\leqslant s\leqslant1\)，所以 \(2-s>0\)，原函数可化为 \(f(x)=\dfrac{2+\sin2x}{4}\). 因此函数的最小正周期为 \(\dfrac{2\pi}{2}=\pi\). 又 \(-1\leqslant\sin2x\leqslant1\)，所以 \(\dfrac14\leqslant f(x)\leqslant\dfrac34\)，最大值为 \(\dfrac34\)，最小值为 \(\dfrac14\).
\end{solution}

\begin{problem}
已知 \(f(x)=ax^3+3x^2-x+1\) 在 \(\R\) 上是减函数，求 \(a\) 的取值范围.
\end{problem}

\begin{solution}
函数在 \(\R\) 上可导. 若函数在 \(\R\) 上为减函数，则对任意 \(x\in\R\) 都有 \(f'(x)\leqslant0\). 反过来，若 \(f'(x)\leqslant0\) 对所有 \(x\in\R\) 成立，且 \(f'\) 不在任何区间上恒为零，则函数在 \(\R\) 上严格递减；本题的 \(f'\) 是二次多项式，满足条件时不会在任何区间上恒为零. 计算得 \(f'(x)=3ax^2+6x-1\).

若 \(a>0\)，则当 \(|x|\) 足够大时 \(3ax^2+6x-1>0\)；若 \(a=0\)，则 \(f'(x)=6x-1\) 不可能对所有实数不大于零. 所以必须有 \(a<0\). 此时 \(f'(x)\) 是开口向下的二次函数，最大值在 \(x=-\dfrac1a\) 处取得，且 \(f'\left(-\dfrac1a\right)=-\dfrac3a-1\). 因而 \(f'(x)\leqslant0\) 对一切 \(x\) 成立，当且仅当 \(-\dfrac3a-1\leqslant0\). 由于 \(a<0\)，该不等式等价于 \(a\leqslant-3\). 当 \(a<-3\) 时，\(f'(x)<0\) 对所有 \(x\in\R\) 成立；当 \(a=-3\) 时，\(f'(x)=-9x^2+6x-1=-(3x-1)^2\leqslant0\)，且它只在 \(x=1/3\) 处为零，因此任意 \(x_1<x_2\) 都有 \(f(x_2)-f(x_1)=\displaystyle\int_{x_1}^{x_2}f'(x)\,\mathrm{d}x<0\). 所以 \(a\leqslant-3\) 时函数确为严格减函数，故 \(a\) 的取值范围为 \((-\infty,-3]\).
\end{solution}

\begin{problem}
从 \(10\) 位同学（其中 \(6\) 女，\(4\) 男）中随机选出 \(3\) 位参加测验. 每位女同学能通过测验的概率均为 \(\dfrac45\)，每位男同学能通过测验的概率均为 \(\dfrac35\)，试求：

\begin{enumerate}
\item 选出的 \(3\) 位同学中，至少有一位男同学的概率；
\item \(10\) 位同学中的女同学甲和男同学乙同时被选中且通过测验的概率.
\end{enumerate}
\end{problem}

\begin{solution}
\begin{enumerate}
\item 从 \(10\) 位同学中等可能选出 \(3\) 位，共有 \(\mathrm C_{10}^{3}\) 种选法. “至少有一位男同学”的对立事件是三位均为女同学，共有 \(\mathrm C_6^3\) 种选法. 所以所求概率为 \(1-\dfrac{\mathrm C_6^3}{\mathrm C_{10}^{3}}=1-\dfrac{20}{120}=\dfrac56\).
\item 同时选中甲、乙时，第三位从其余 \(8\) 位同学中任选，故甲、乙同时被选中的概率为 \(\dfrac{\mathrm C_8^1}{\mathrm C_{10}^{3}}=\dfrac1{15}\). 在已选中甲、乙的条件下，两人同时通过测验的概率为 \(\dfrac45\cdot\dfrac35=\dfrac{12}{25}\). 选人和测验结果相互独立，因此所求概率为 \(\dfrac1{15}\cdot\dfrac45\cdot\dfrac35=\dfrac4{125}\).
\end{enumerate}
\end{solution}

\begin{problem}
如图，已知四棱锥 \(P-ABCD\)，\(PB\perp AD\)，侧面 \(PAD\) 为边长等于 \(2\) 的正三角形，底面 \(ABCD\) 为菱形，侧面 \(PAD\) 与底面 \(ABCD\) 所成的二面角为 \(120^{\circ}\).

\begin{enumerate}
\item 求点 \(P\) 到平面 \(ABCD\) 的距离；
\item 求面 \(APB\) 与面 \(CPB\) 所成二面角的大小.
\end{enumerate}

\begin{center}
\bitmapfigure[max width=0.42\linewidth]{img/2004/national_paper_1_liberal/q21_fig1.png}
\end{center}
\end{problem}

\begin{solution}
\begin{enumerate}
\item 设 \(O\) 为点 \(P\) 在平面 \(ABCD\) 上的射影，连接 \(OB\)，令 \(OB\cap AD=E\). 因为 \(PO\perp AD\)，题设又有 \(PB\perp AD\)，且 \(PO\)，\(PB\) 是平面 \(POB\) 内相交的两条直线，所以 \(AD\perp\) 平面 \(POB\)，从而 \(AD\perp PE\). 在正三角形 \(PAD\) 中，过顶点 \(P\) 的高也是中线，所以 \(E\) 是 \(AD\) 的中点，\(AE=DE=1\)，且 \(PE=\sqrt{2^2-1^2}=\sqrt3\).

二面角 \(P-AD-B\) 的平面角为 \(\angle PEB=120^{\circ}\). 因为 \(E\)，\(O\)，\(B\) 共线，故 \(\angle PEO=60^{\circ}\). 在直角三角形 \(POE\) 中，\(PO=PE\sin60^{\circ}=\sqrt3\cdot\dfrac{\sqrt3}{2}=\dfrac32\). 因而点 \(P\) 到平面 \(ABCD\) 的距离为 \(\dfrac32\).
\item 以 \(E\) 为原点，以 \(EA\)，\(EB\) 和垂直于底面且与 \(OP\) 同向的方向分别为 \(x\)，\(y\)，\(z\) 轴正方向. 因为 \(EA=ED=1\)，可取 \(A(1,0,0)\)，\(D(-1,0,0)\). 由 \(EA\perp EB\) 且 \(AB=AD=2\)，可得 \(B(0,\sqrt3,0)\)，\(C(-2,\sqrt3,0)\). 又由 \(PE=\sqrt3\)，\(PO=\dfrac32\) 及 \(\angle PEB=120^{\circ}\)，可取 \(P\left(0,-\dfrac{\sqrt3}{2},\dfrac32\right)\).
于是 \(\overrightarrow{AB}=(-1,\sqrt3,0)\)，\(\overrightarrow{BP}=\left(0,-\dfrac{3\sqrt3}{2},\dfrac32\right)\)，\(\overrightarrow{BC}=(-2,0,0)\).

设平面 \(APB\) 的一个法向量为 \(\bs{m}=(m_1,m_2,m_3)\). 由 \(\bs{m}\cdot\overrightarrow{AB}=0\) 和 \(\bs{m}\cdot\overrightarrow{BP}=0\)，即 \(-m_1+\sqrt3m_2=0\)，\(-\dfrac{3\sqrt3}{2}m_2+\dfrac32m_3=0\)，取 \(m_2=1\)，得 \(\bs{m}=(\sqrt3,1,\sqrt3)\). 设平面 \(CPB\) 的一个法向量为 \(\bs{n}=(n_1,n_2,n_3)\). 由 \(\bs{n}\cdot\overrightarrow{BC}=0\) 和 \(\bs{n}\cdot\overrightarrow{BP}=0\)，即 \(-2n_1=0\)，\(-\dfrac{3\sqrt3}{2}n_2+\dfrac32n_3=0\)，取 \(n_2=1\)，得 \(\bs{n}=(0,1,\sqrt3)\).

两法向量的夹角满足
\[
\cos\angle(\bs{m},\bs{n})=\frac{\bs{m}\cdot\bs{n}}{|\bs{m}|\,|\bs{n}|}=\frac4{\sqrt7\cdot2}=\frac{2\sqrt7}{7}
\]
上式给出两法向量的锐角夹角，而题目所求的两个面的内二面角为其补角，所以所求二面角为 \(\pi-\arccos\dfrac{2\sqrt7}{7}\).
\end{enumerate}
\end{solution}

\begin{problem}
设双曲线 \(C:\dfrac{x^2}{a^2}-y^2=1\)（\(a>0\)）与直线 \(l:x+y=1\) 相交于两个不同的点 \(A\)，\(B\).

\begin{enumerate}
\item 求双曲线 \(C\) 的离心率 \(e\) 的取值范围；
\item 设直线 \(l\) 与 \(y\) 轴的交点为 \(P\)，且 \(\overrightarrow{PA}=\dfrac5{12}\overrightarrow{PB}\). 求 \(a\) 的值.
\end{enumerate}
\end{problem}

\begin{solution}
\begin{enumerate}
\item 将 \(l\) 的方程 \(y=1-x\) 代入双曲线方程，得 \(\dfrac{x^2}{a^2}-(1-x)^2=1\). 由于 \(a>0\)，乘以 \(a^2\) 并整理得
\[
(1-a^2)x^2+2a^2x-2a^2=0
\]
若 \(a=1\)，上式退化为一次方程，直线至多与双曲线有一个交点，所以必须有 \(a\ne1\). 当 \(a\ne1\) 时，方程有两个不同实根的充要条件为
\[
\Delta=(2a^2)^2-4(1-a^2)(-2a^2)=4a^2(2-a^2)>0
\]
结合 \(a>0\)，得 \(0<a<\sqrt2\)，再排除 \(a=1\)，所以 \(a\in(0,1)\cup(1,\sqrt2)\). 双曲线的离心率为 \(e=\sqrt{1+\dfrac1{a^2}}\). 当 \(0<a<1\) 时，\(e\in(\sqrt2,+\infty)\)；当 \(1<a<\sqrt2\) 时，\(e\in\left(\sqrt{\dfrac32},\sqrt2\right)\). 因此 \(e\in\left(\sqrt{\dfrac32},\sqrt2\right)\cup\left(\sqrt2,+\infty\right)\).
\item 直线 \(l\) 与 \(y\) 轴交于 \(P(0,1)\). 在直线 \(l\) 上用参数 \(t\) 表示点 \(X(t)=(t,1-t)\)，则 \(\overrightarrow{PX(t)}=t(1,-1)\). 设对应于 \(A\)，\(B\) 的参数分别为 \(t_A\)，\(t_B\). 由题设向量关系，得 \(t_A=\dfrac5{12}t_B\). 令 \(t_A=5k\)，\(t_B=12k\)，则二次方程的两个根为 \(5k\)，\(12k\).

由韦达定理，根的和、根的积分别为 \(-\dfrac{2a^2}{1-a^2}\) 和 \(-\dfrac{2a^2}{1-a^2}\)，所以二者相等. 因而 \(17k=60k^2\). 因交点不同，\(k\ne0\)，所以 \(k=\dfrac{17}{60}\). 根的和为 \(17k=\dfrac{289}{60}\)，故
\[
-\frac{2a^2}{1-a^2}=\frac{289}{60}
\]
整理得 \(169a^2=289\). 因 \(a>0\)，所以 \(a=\dfrac{17}{13}\).
\end{enumerate}
\end{solution}
